% Laget automatisk av elevark/lag_elevark.py fra sim/heis.html. Endringer her blir overskrevet.
\documentclass[11pt]{article}
\usepackage{elevark}
\begin{document}
\elevarkhode{Fysikk 1 · Krefter}{Baderomsvekt i heisen}{https://fysikklab.no/sim/heis.html}{Simuleringen}

Du står på en baderomsvekt i en heis. Vekta måler hvor hardt du trykker på den, altså normalkraften. Hva viser den når heisen starter, kjører jevnt og stopper? Og hva hvis kabelen ryker?

\section*{Oppgaver}

Start med en person på 60~kg, tur opp, akselerasjon 1,5~m/s² og toppfart 2,0~m/s.

\begin{oppgaveliste}
\oppgave \fase{gjett}{Gjett}Viser vekta mer eller mindre enn 60~kg når heisen starter oppover? Når den kjører med konstant fart? Når den bremser?
\svar
\oppgave \fase{forklar}{Forklar}Tegn kreftene på personen. Bruk Newtons 2. lov til å vise at $N = m(g + a)$ når positiv retning er oppover.
\svar
\oppgave \fase{test}{Regn}Regn ut normalkraften og hva vekta viser i hver av de tre delene av turen opp. Sjekk med simuleringen.
\svar
\oppgave \fase{gjett}{Gjett}Hva viser vekta når heisen starter nedover og når den bremser på vei ned? Test.
\svar
\oppgave \fase{test}{Regn}I forsøket med baderomsvekt i heisen veier en elev 58~kg i ro. Når heisen starter, viser vekta 65~kg. Hvor stor er akselerasjonen, og hvilken vei går heisen? Still inn simuleringen slik at du får det samme.
\svar
\oppgave \fase{forklar}{Forklar}Velg «Kabelen ryker». Hva viser vekta? Hvorfor sier vi at astronauter i en romstasjon er vektløse, selv om tyngdekraften fra Jorda virker på dem?
\svar
\oppgave \fase{utfordring}{Utfordring}Kan du ut fra grafen for vekta finne ut om heisen gikk opp eller ned? Kan du finne ut hvor fort den kjørte?
\svar
\end{oppgaveliste}

\end{document}
